\section{Solutions for Chapter~\ref{ch:scaling}}

\paragraph{Exercise~\ref{ex:scaling-residual}.}
\phantomsection
\label{sol:ex:scaling-residual}
Expand $u=\sum_nc_n|n\rangle$. Then
$\|r\|^2=\sum_n|c_n|^2(E_n-E)^2$. This is at least the squared
distance from $E$ to the spectrum, since the weights sum to one.
If $\Pi$ projects onto the target cluster and all other eigenvalues
are at least $\gamma$ from $E$, the same sum is at least
$\gamma^2\|(1-\Pi)u\|^2$. Hence the unwanted squared weight is
at most $\|r\|^2/\gamma^2$. The value of $\gamma$ is a separate
spectral input.



\paragraph{Exercise~\ref{ex:scaling-stress}.}
\phantomsection
\label{sol:ex:scaling-stress}
For the vacuum, $L_n|0\rangle=0$ for $n\geq-1$ and
$L_0|0\rangle=0$. Therefore
\[
 \|L_{-n}|0\rangle\|^2
 =\langle0|[L_n,L_{-n}]|0\rangle
 =\frac c{12}n(n^2-1).
\]
At $n=2$ this is $c/2$. The source amplitude squared is
$(2\pi v/L)^2c/2$, whereas the stress gap squared is
$(4\pi v/L)^2$. Their ratio is $c/8$, proving the estimator.
The calculation assumes the source has the displayed Virasoro
normalization; it does not establish the lattice-to-CFT map.



\paragraph{Exercise~\ref{ex:scaling-improvement}.}
\phantomsection
\label{sol:ex:scaling-improvement}
Put $k=2\pi m/L$ and use periodic indices. In the first sum
below, relabel $i=j+1$:
\[
 \sum_j e^{-ikj}(B_{j+1}-B_j)
 =e^{ik}\sum_i e^{-iki}B_i-\sum_j e^{-ikj}B_j.
\]
The phase is periodic because $e^{ikL}=1$, so the difference
is $(e^{ik}-1)B_m$. It vanishes identically for $m=0$, as
required for the unchanged total Hamiltonian.

For fixed nonzero $m$, matrix elements of a scaling field
$B_j$ of dimension $\Delta_B$ between normalized low-energy
states scale as $L^{-\Delta_B}$. Summing its Fourier component
over $L$ sites gives $O(L^{1-\Delta_B})$ when the relevant
continuum matrix element is nonzero. Since
$e^{ik}-1=2\pi i m/L+O(L^{-2})$, the improvement contributes
$O(L^{-\Delta_B})$. A dimension-one current can therefore
contribute at the same $L^{-1}$ order as the extensive stress
Fourier mode. A field of dimension greater than one is
subleading under this scaling assumption. Selection rules may
make a particular matrix element vanish; dimension counting
alone does not guarantee a nonzero coefficient.



\paragraph{Exercise~\ref{ex:scaling-casimir}.}
\phantomsection
\label{sol:ex:scaling-casimir}
Write
$E_0(L)/L=e_\infty-A/L^2+r(L)$, with $A=\pi vc/6$.
Subtracting the two sizes eliminates $e_\infty$ and gives
\begin{align*}
 (vc)_{\mathrm{two\ sizes}}
 &=-\frac6\pi
 \frac{E_0(L_1)/L_1-E_0(L_2)/L_2}
 {L_1^{-2}-L_2^{-2}},\\
 &=vc-\frac6\pi\frac{r(L_1)-r(L_2)}
 {L_1^{-2}-L_2^{-2}}.
\end{align*}
Thus even small per-site remainders can matter when divided by
the difference of two small inverse-square lengths. Two exact
energies solve for two leading coefficients exactly, but do not
estimate this remainder.



\paragraph{Exercise~\ref{ex:scaling-fraction}.}
\phantomsection
\label{sol:ex:scaling-fraction}
If $\phi$ is scalar of dimension $\Delta_\phi$, then
$L_{-1}\phi$ has weights
$(\Delta_\phi/2+1,\Delta_\phi/2)$ and spin one. A holomorphic
current has weights $(1,0)$. The former has nonzero right weight
whenever $\Delta_\phi>0$. A momentum-one level near
$\Delta_\phi+1$ is therefore consistent with an ordinary
descendant and is not evidence for a current. It could also mix
with an unrelated finite-size primary.



\paragraph{Exercise~\ref{ex:scaling-modular}.}
\phantomsection
\label{sol:ex:scaling-modular}
Let $D=\beta\partial_\beta$, $a=\pi c/(6\beta_*)$, and
$\mu_\beta=\langle g\rangle_\beta$. Directly differentiating
$Z=e^{\beta a}\sum_j e^{-\beta g_j}$ gives
\[
 \frac{D Z}{Z}=\beta(a-\mu_\beta).
\]
Its vanishing at $\beta=\beta_*$ gives $a=\mu_{\beta_*}=\mu$,
or $c=6\beta_*\mu/\pi$. At that point,
$Z^{(j)}/Z=(-1)^jm_j$ for $j\geq2$, where the derivatives
are in $\beta$ and the moments are centered. Repeated use of
$D\beta^j=\beta^j(D+j)$ gives
\begin{align*}
 D^3&=\beta\partial_\beta+3\beta^2\partial_\beta^2+
       \beta^3\partial_\beta^3,\\
 D^5&=\beta\partial_\beta+15\beta^2\partial_\beta^2+
       25\beta^3\partial_\beta^3\\
 &\hspace{2em}+10\beta^4\partial_\beta^4+
       \beta^5\partial_\beta^5.
\end{align*}
The first-derivative term vanishes. The third condition is
$3\beta_*^2m_2-\beta_*^3m_3=0$, hence $\beta_*m_3=3m_2$.
The fifth condition is
\[
 15\beta_*^2m_2-25\beta_*^3m_3
       +10\beta_*^4m_4-\beta_*^5m_5=0.
\]
These are derivatives of $Z$ itself, so the $m_j$ are central
moments; replacing them by cumulants would require differentiating
$\log Z$ instead. The odd
derivatives vanish because modular $S$ acts by $t\mapsto-t$
in $\beta=\beta_*e^t$. A charge projection inserts a symmetry
operator into the torus trace. Exchanging cycles generally
turns that insertion into a spatial twist; the projected
partition function alone is therefore not $S$-invariant.


